\documentclass[10pt,a4paper]{amsart}
\usepackage{amsmath,amssymb,amsthm,graphicx}
\usepackage{verbatim}
\newenvironment{reference}{\comment}{\endcomment}
%\newenvironment{reference}{}{}
\newenvironment{slogan}{\comment}{\endcomment}
\newenvironment{history}{\comment}{\endcomment}
% For commutative diagrams we use Xy-pic
\usepackage[all]{xy}

% We use 2cell for 2-commutative diagrams.
\xyoption{2cell}
% Theorem environments.
%
\theoremstyle{plain}
\newtheorem{theorem}[subsection]{Theorem}
\newtheorem{proposition}[subsection]{Proposition}
\newtheorem{lemma}[subsection]{Lemma}

\theoremstyle{definition}
\newtheorem{definition}[subsection]{Definition}
\newtheorem{example}[subsection]{Example}
\newtheorem{exercise}[subsection]{Exercise}
\newtheorem{situation}[subsection]{Situation}

\theoremstyle{remark}
\newtheorem{remark}[subsection]{Remark}
\newtheorem{remarks}[subsection]{Remarks}

\numberwithin{equation}{subsection}

% Macros
%
\def\lim{\mathop{\mathrm{lim}}\nolimits}
\def\colim{\mathop{\mathrm{colim}}\nolimits}
\def\Spec{\mathop{\mathrm{Spec}}}
\def\Hom{\mathop{\mathrm{Hom}}\nolimits}
\def\Ext{\mathop{\mathrm{Ext}}\nolimits}
\def\SheafHom{\mathop{\mathcal{H}\!\mathit{om}}\nolimits}
\def\SheafExt{\mathop{\mathcal{E}\!\mathit{xt}}\nolimits}
\def\Sch{\mathit{Sch}}
\def\Mor{\mathop{\mathrm{Mor}}\nolimits}
\def\Ob{\mathop{\mathrm{Ob}}\nolimits}
\def\Sh{\mathop{\mathit{Sh}}\nolimits}
\def\NL{\mathop{N\!L}\nolimits}
\def\CH{\mathop{\mathrm{CH}}\nolimits}
\def\proetale{{pro\text{-}\acute{e}tale}}
\def\etale{{\acute{e}tale}}
\def\QCoh{\mathit{QCoh}}
\def\Ker{\mathop{\mathrm{Ker}}}
\def\Im{\mathop{\mathrm{Im}}}
\def\Coker{\mathop{\mathrm{Coker}}}
\def\Coim{\mathop{\mathrm{Coim}}}

% Boxtimes
%
\DeclareMathSymbol{\boxtimes}{\mathbin}{AMSa}{"02}

%
% Macros for moduli stacks/spaces
%
\def\QCohstack{\mathcal{QC}\!\mathit{oh}}
\def\Cohstack{\mathcal{C}\!\mathit{oh}}
\def\Spacesstack{\mathcal{S}\!\mathit{paces}}
\def\Quotfunctor{\mathrm{Quot}}
\def\Hilbfunctor{\mathrm{Hilb}}
\def\Curvesstack{\mathcal{C}\!\mathit{urves}}
\def\Polarizedstack{\mathcal{P}\!\mathit{olarized}}
\def\Complexesstack{\mathcal{C}\!\mathit{omplexes}}
% \Pic is the operator that assigns to X its picard group, usage \Pic(X)
% \Picardstack_{X/B} denotes the Picard stack of X over B
% \Picardfunctor_{X/B} denotes the Picard functor of X over B
\def\Pic{\mathop{\mathrm{Pic}}\nolimits}
\def\Picardstack{\mathcal{P}\!\mathit{ic}}
\def\Picardfunctor{\mathrm{Pic}}
\def\Deformationcategory{\mathcal{D}\!\mathit{ef}}
\usepackage{cleveref}
\crefname{originalreference}{original reference}{original references}
\Crefname{originalreference}{Original reference}{Original references}
\makeatletter
\newlabel{algebra-lemma-hom-from-finitely-presented}{{algebra-lemma-hom-from-finitely-presented}{source}{}{sourcepointer-1}{}}
\newlabel{algebra-lemma-hom-from-finitely-presented@cref}{{[originalreference][0][]algebra-lemma-hom-from-finitely-presented}{[1][0][]source}{}{}{}}
\makeatother
\crefname{originalreference}{original reference}{original references}
\Crefname{originalreference}{Original reference}{Original references}
\makeatletter
\newlabel{injectives-lemma-criterion-baer}{{injectives-lemma-criterion-baer}{source}{}{sourcepointer-2}{}}
\newlabel{injectives-lemma-criterion-baer@cref}{{[originalreference][0][]injectives-lemma-criterion-baer}{[1][0][]source}{}{}{}}
\makeatother
\crefname{originallemma}{original lemma}{original lemmas}
\Crefname{originallemma}{Original lemma}{Original lemmas}
\makeatletter
\newlabel{lemma-injective-epimorphism}{{lemma-injective-epimorphism}{source}{}{sourcepointer-3}{}}
\newlabel{lemma-injective-epimorphism@cref}{{[originallemma][0][]lemma-injective-epimorphism}{[1][0][]source}{}{}{}}
\makeatother
\title{Localizing an injective module over a Noetherian ring stays injective}
\author{The Stacks Project Authors}
\date{stacks-project:dualizing.tex:local-snapshot-20261012; GNU FDL 1.2 or later}
\begin{document}
\maketitle
\noindent\textit{Editorial note.} All delivered mathematical text is the author's own, copied from the pinned Stacks Project chapter source sources/stacks-dualizing/pinned-original-source.tex; no statement or proof step was rewritten, reordered or removed.\par\medskip
\noindent\textit{Editorial note.} Preamble provenance: the Stacks Project keeps one single shared preamble for every chapter; this card takes it directly from the staged sources/stacks-dualizing/preamble.tex as three verbatim spans (lines 8-12, 14-18 and 154-217), so no derived file was written and nothing was added or redefined. The wrapper class and the base packages are supplied by the assembler.\par\medskip
\noindent\textit{Editorial note.} Bibliography provenance: neither the counted statement nor its proof contains a citation (this lane rescans the whole span before writing the card), so the card carries no bibliography and no entry was reconstructed.\par\medskip
\noindent\textit{Editorial note.} Selection provenance: this card was opened under the file-based control contract (CONTROL.md), with the mechanical pre-card gate run on its statement and proof interval before the card was written.\par\medskip
\noindent\textit{Allowed standard prerequisites (editorial).}\begin{itemize}
\item The reduction of injectivity over a ring epimorphism to injectivity over the source ring. Invoked at the start of the counted proof (source lines 534-536). Reduces the claim about the localized module to the claim that it is injective as a module over the original ring.
\item Baer's criterion for injectivity of a module. Invoked inside the counted proof (source line 537). Reduces injectivity to extending maps out of ideals.
\item The lifting of a map out of a finitely presented module after inverting an element of the multiplicative subset. Invoked inside the counted proof (source lines 540-543). Produces the element of the multiplicative subset and the map into the original module whose composition is the given map.
\end{itemize}\medskip
\ifdefined\hypertarget\hypertarget{sourcepointer-1}{}\fi
\noindent\textit{Source reference.} Original reference algebra-lemma-hom-from-finitely-presented: sources/stacks-dualizing/pinned-original-source.tex, source lines 543--543. Navigation pointer only: the counted text refers to the Stacks item algebra-lemma-hom-from-finitely-presented, which belongs to another chapter of the project and is not part of this card. The pointer carries the target's own label as its number because this chapter's staged snapshot does not contain the other chapter's tag table; its content is not reproduced here.\par\medskip
\ifdefined\hypertarget\hypertarget{sourcepointer-2}{}\fi
\noindent\textit{Source reference.} Original reference injectives-lemma-criterion-baer: sources/stacks-dualizing/pinned-original-source.tex, source lines 517--517. Navigation pointer only: the counted text refers to the Stacks item injectives-lemma-criterion-baer, which belongs to another chapter of the project and is not part of this card. The pointer carries the target's own label as its number because this chapter's staged snapshot does not contain the other chapter's tag table; its content is not reproduced here.\par\medskip
\ifdefined\hypertarget\hypertarget{sourcepointer-3}{}\fi
\noindent\textit{Source reference.} Original lemma lemma-injective-epimorphism: sources/stacks-dualizing/pinned-original-source.tex, source lines 422--422. Navigation pointer only: the counted text refers to the Stacks item lemma-injective-epimorphism, which is an existing labeled result of this same chapter and is not part of this card. The pointer carries the target's own label as its number and its definition line as its span; its content is not reproduced here.\par\medskip
\begin{lemma}
\label{lemma-localization-injective-modules}
Let $R$ be a Noetherian ring. Let $S \subset R$ be a multiplicative
subset. If $E$ is an injective $R$-module, then $S^{-1}E$ is an
injective $S^{-1}R$-module.
\end{lemma}

\begin{proof}
Since $R \to S^{-1}R$ is an epimorphism of rings, it suffices
to show that $S^{-1}E$ is injective as an $R$-module, see
Lemma \ref{lemma-injective-epimorphism}.
To show this we use Injectives, Lemma \ref{injectives-lemma-criterion-baer}.
Thus let $I \subset R$ be an ideal and let
$\varphi : I \to S^{-1} E$ be an $R$-module map.
As $I$ is a finitely presented $R$-module (because $R$ is Noetherian)
we can find an $f \in S$ and an $R$-module map $I \to E$
such that $f\varphi$ is the composition $I \to E \to S^{-1}E$
(Algebra, Lemma \ref{algebra-lemma-hom-from-finitely-presented}).
Then we can extend $I \to E$ to a homomorphism $R \to E$.
Then the composition
$$
R \to E \to S^{-1}E \xrightarrow{f^{-1}} S^{-1}E
$$
is the desired extension of $\varphi$ to $R$.
\end{proof}

\end{document}
